\section{R ile çözümlü uygulama}
\label{sec:r-b07}

\textbf{Soru.} Aşağıdaki on ikili yanıt için oran tahminini ve yaklaşık
standart hatayı bulun. Ardından gerçek oranı 0,40 olan Bernoulli evreninde
$n=50$ için tahmin edicinin yanlılığını ve MSE'sini benzetimle inceleyin.

\begin{lstlisting}[style=prismR]
yanit <- c(1, 0, 1, 1, 0, 1, 1, 0, 1, 0)
oran <- mean(yanit)
c(oran = oran,
  standart_hata = sqrt(oran * (1 - oran) / length(yanit)))
set.seed(2026)
gercek_oran <- 0.40
hacim <- 50
tahminler <- rbinom(10000, size = hacim,
                    prob = gercek_oran) / hacim
c(merkez = mean(tahminler),
  yanlilik = mean(tahminler) - gercek_oran,
  ampirik_se = sd(tahminler),
  mse = mean((tahminler - gercek_oran)^2))
c(kuramsal_se = sqrt(0.40 * 0.60 / hacim),
  kuramsal_mse = 0.40 * 0.60 / hacim)
hacimler <- c(25, 400)
sqrt(0.60 * 0.40 / hacimler)
\end{lstlisting}

\begin{outputguidebox}
\textbf{Çözüm ve kontrol değerleri.} On yanıtta altı olumlu sonuç vardır:
$\hat p=0{,}60$ ve yaklaşık $\widehat{\SE}=0{,}1549$.
Benzetimin merkezi yaklaşık 0,40, yanlılığı yaklaşık sıfır olmalıdır.
Kuramsal standart hata 0,06928, MSE 0,0048'dir. $\hat p=0{,}60$ için
$n=25$ ve $n=400$ kullanılırsa yaklaşık standart hatalar 0,09798 ve 0,02449 olur.
\end{outputguidebox}

\textbf{Yorum.} İlk hesap bilinmeyen oran yerine örneklem oranını koyar;
benzetimin kuramsal kontrolü ise bilinen 0,40 oranını kullanır. Sonlu sayıdaki
tekrarda görülen küçük ampirik yanlılık, tahmin edicinin kuramsal olarak
yanlı olduğunu kanıtlamaz. On gözlem için normal yaklaşım temkinle kullanılmalıdır.

\textbf{Pekiştirme ve yanıt.} Bernoulli örneklem oranı yansız olduğundan
kuramsal MSE varyansa eşittir: $p(1-p)/n=0{,}0048$.